\textbf{Phần 1. Trắc nghiệm lựa chọn}
\setcounterpageref{ex}{0}

\begin{ex} %[2H11N56]
Đơn vị tiền tệ của nước Việt Nam là
\choice
{Đô-la}
{Yên}
{\True Đồng}
{Ơ-rô}
\loigiai{
Đơn vị tiền tệ là Đồng.\\
\textbf{Đáp án của câu hỏi là C}
}
\end{ex}

\begin{ex} %[2H11N56]
Tờ tiền nào sau đây có mệnh giá lớn nhất?
\choice
{$100$ đồng}
{\True $1\,000$ đồng}
{$200$ đồng}
{$500$ đồng}
\loigiai{
Tờ 1 000 đồng lớn nhất.\\
\textbf{Đáp án của câu hỏi là B}
}
\end{ex}

\begin{ex} %[2H11H56]
2 tờ 500 đồng có giá trị bằng
\choice
{\True $1\,000$ đồng}
{$500$ đồng}
{$700$ đồng}
{$2\,000$ đồng}
\loigiai{
$500 + 500 = 1000$ đồng.\\
\textbf{Đáp án của câu hỏi là A}
}
\end{ex}

\begin{ex} %[2H11H56]
5 tờ 200 đồng có tổng giá trị là
\choice
{$800$ đồng}
{\True $1\,000$ đồng}
{$700$ đồng}
{$500$ đồng}
\loigiai{
$200 \times 5 = 1000$ đồng.\\
\textbf{Đáp án của câu hỏi là B}
}
\end{ex}

\begin{ex} %[2H11H56]
Mai có 1 tờ 500 đồng và 2 tờ 200 đồng. Mai có tất cả số tiền là
\choice
{\True $900$ đồng}
{$800$ đồng}
{$1000$ đồng}
{$700$ đồng}
\loigiai{
$500 + 200 + 200 = 900$ đồng.\\
\textbf{Đáp án của câu hỏi là A}
}
\end{ex}

\begin{ex} %[2H11H56]
Một chiếc kẹo mút giá 500 đồng, một cái bánh que giá 300 đồng. Mua cả hai thứ hết số tiền là
\choice
{$700$ đồng}
{$900$ đồng}
{$1000$ đồng}
{\True $800$ đồng}
\loigiai{
$500 + 300 = 800$ đồng.\\
\textbf{Đáp án của câu hỏi là D}
}
\end{ex}

\begin{ex} %[2H11V56]
Nam có 1 000 đồng, Nam mua một cục tẩy hết 600 đồng. Nam còn lại số tiền là
\choice
{$300$ đồng}
{$500$ đồng}
{\True $400$ đồng}
{$200$ đồng}
\loigiai{
$1000 - 600 = 400$ đồng.\\
\textbf{Đáp án của câu hỏi là C}
}
\end{ex}

\begin{ex} %[2H11V56]
Để có đủ 1 000 đồng từ các tờ 200 đồng, cần mấy tờ 200 đồng?
\choice
{$4$ tờ}
{$6$ tờ}
{$10$ tờ}
{\True $5$ tờ}
\loigiai{
$1000 : 200 = 5$ tờ.\\
\textbf{Đáp án của câu hỏi là D}
}
\end{ex}

\begin{ex} %[2H11V56]
Bình có 3 tờ 200 đồng và 1 tờ 100 đồng. Bình có tất cả số tiền là
\choice
{\True $700$ đồng}
{$600$ đồng}
{$800$ đồng}
{$900$ đồng}
\loigiai{
$200 \times 3 + 100 = 700$ đồng.\\
\textbf{Đáp án của câu hỏi là A}
}
\end{ex}

\begin{ex} %[2H11V56]
Một chiếc bút chì giá 800 đồng, chiếc thước kẻ rẻ hơn chiếc bút chì 200 đồng. Chiếc thước kẻ giá
\choice
{$500$ đồng}
{\True $600$ đồng}
{$700$ đồng}
{$1000$ đồng}
\loigiai{
$800 - 200 = 600$ đồng.\\
\textbf{Đáp án của câu hỏi là B}
}
\end{ex}

\begin{ex} %[2H11V56]
Lan có 1 000 đồng, Lan mua 2 chiếc kẹo giá 400 đồng một chiếc. Lan còn lại số tiền là
\choice
{$100$ đồng}
{$300$ đồng}
{\True $200$ đồng}
{$400$ đồng}
\loigiai{
Mua 2 kẹo: $400 \times 2 = 800$ đồng.\\
Còn lại: $1000 - 800 = 200$ đồng.\\
\textbf{Đáp án của câu hỏi là C}
}
\end{ex}

\begin{ex} %[2H11C56]
Mẹ đưa cho cô bán hàng tờ 1 000 đồng để mua một mớ rau giá 700 đồng. Cô bán hàng trả lại tiền thừa bằng các tờ tiền nào sau đây?
\choice
{1 tờ 200 đồng}
{2 tờ 200 đồng}
{3 tờ 200 đồng}
{\True 1 tờ 200 đồng và 1 tờ 100 đồng}
\loigiai{
Tiền thừa: $1000 - 700 = 300$ đồng = $200 + 100$ đồng.\\
\textbf{Đáp án của câu hỏi là D}
}
\end{ex}

\textbf{Phần 2. Trắc nghiệm đúng sai}
\setcounterpageref{ex}{0}

\begin{ex} %[2H11H56]
Xét tính đúng sai về mệnh giá tiền Việt Nam:
\choiceTF
{\True 2 tờ 500 đồng có giá trị bằng 1 000 đồng.}
{\True 5 tờ 100 đồng có giá trị bằng 500 đồng.}
{1 tờ 1 000 đồng có giá trị bằng 3 tờ 200 đồng.}
{\True 1 tờ 1 000 đồng có giá trị bằng 5 tờ 200 đồng.}
\loigiai{
\begin{itemize}
    \item $500 \times 2 = 1000$ đồng: Đúng.
    \item $100 \times 5 = 500$ đồng: Đúng.
    \item 1000 đồng = 5 tờ 200 đồng (không phải 3 tờ). Ý 3 Sai, Ý 4 Đúng.
\end{itemize}
\begin{itemize}
    \item a) \textbf{Đúng}.
    \item b) \textbf{Đúng}.
    \item c) \textbf{Sai}.
    \item d) \textbf{Đúng}.
\end{itemize}
}
\end{ex}

\begin{ex} %[2H11V56]
Hà có 1 000 đồng, Hà mua 1 gói kẹo giá 500 đồng và 1 phong bì giá 200 đồng. Xét tính đúng sai:
\choiceTF
{\True Tổng số tiền Hà đã mua là 700 đồng ($500 + 200 = 700$).}
{\True Số tiền Hà còn lại là 300 đồng ($1000 - 700 = 300$).}
{Hà đã tiêu hết toàn bộ 1 000 đồng.}
{\True Hà vẫn còn thừa 300 đồng.}
\loigiai{
\begin{itemize}
    \item Mua hết 700 đồng: Đúng.
    \item Còn lại 300 đồng: Đúng.
    \item Chưa tiêu hết. Ý 3 Sai, Ý 4 Đúng.
\end{itemize}
\begin{itemize}
    \item a) \textbf{Đúng}.
    \item b) \textbf{Đúng}.
    \item c) \textbf{Sai}.
    \item d) \textbf{Đúng}.
\end{itemize}
}
\end{ex}

\textbf{Phần 3. Trả lời ngắn}
\setcounterpageref{ex}{0}

\begin{ex} %[2H11H56]
2 tờ 500 đồng đổi được bao nhiêu đồng?
\shortans{1000}
\loigiai{
\textbf{Đáp số:} $1000$\\
1000 đồng.
}
\end{ex}

\begin{ex} %[2H11H56]
Có 4 tờ 200 đồng, hỏi có tất cả bao nhiêu đồng?
\shortans{800}
\loigiai{
\textbf{Đáp số:} $800$\\
$200 \times 4 = 800$ đồng.
}
\end{ex}

\begin{ex} %[2H11V56]
Lan có 1 000 đồng, mua kem hết 700 đồng. Hỏi Lan còn lại bao nhiêu đồng?
\shortans{300}
\loigiai{
\textbf{Đáp số:} $300$\\
$1000 - 700 = 300$ đồng.
}
\end{ex}

\begin{ex} %[2H11V56]
Một cái nhãn vở giá 200 đồng. Hỏi mua 4 cái nhãn vở hết bao nhiêu đồng?
\shortans{800}
\loigiai{
\textbf{Đáp số:} $800$\\
$200 \times 4 = 800$ đồng.
}
\end{ex}

\begin{ex} %[2H11V56]
Bình có 1 tờ 500 đồng và 2 tờ 200 đồng. Hỏi Bình có tất cả bao nhiêu đồng?
\shortans{900}
\loigiai{
\textbf{Đáp số:} $900$\\
$500 + 400 = 900$ đồng.
}
\end{ex}

\begin{ex} %[2H11C56]
Cần bao nhiêu tờ 200 đồng để đổi được 1 tờ 1 000 đồng?
\shortans{5}
\loigiai{
\textbf{Đáp số:} $5$\\
$1000 : 200 = 5$ tờ.
}
\end{ex}
